\subsection{Actual doubled corrections across the full class-two source}

\begin{statement}{D.54}{Proved}{Integrated even corrections, the reduced arithmetic boundary map, and the complete finite comparison}
Keep the actual arithmetic source and the fixed normalized
conjugators of D.50. Write
\[
 \begin{gathered}
 R_{50}=\mathfrak R_h\cap\mathfrak K_\infty,\qquad
 L_2=\overline{\Gamma_3(K_3)(K_3')^2K_3^4},\\
 \mathcal T_2=\widehat B/L_2,\qquad
 D=\ker(\mathcal T_2\to G_3),\qquad
 C=D'=\bigwedge\nolimits^2_{\mathbb F_2}V,\\
 Z=W\wedge V,\qquad Z^0=Z^A,\qquad
 Z^-=(1-e_A)Z .
 \end{gathered}
 \tag{D.54.1}
\]
D.51 proved that the exact image of \(R_{50}\) in
\(\mathcal T_2\) is \(Z\).
Consequently the full joint source and its two
specified projections are
\[
 \begin{gathered}
 \mathcal X_{54}
   =\widehat B/(R_{50}\cap L_2)
   =\mathcal Z_h\times_{\mathcal T_2/Z}\mathcal T_2,\\
 \ker(\mathcal X_{54}\to\mathcal Z_h)=Z,\qquad
 \ker(\mathcal X_{54}\to\mathcal T_2)
       =L_2R_{50}/R_{50}.
 \end{gathered}
 \tag{D.54.2}
\]
The map to D.51's joint source has kernel \(Z^0\).

There is an actual peripheral automorphism of
\(\mathcal X_{54}\) for every doubled integral
pair \(2(b,c)\), \((b,c)\in\mathscr P\) of
D.53.10. These automorphisms act identically
on \(\mathcal W_h\), fix \(t=sr\), and have
complete homology action
\[
                  1+2D_{b,c}\quad\hbox{on }U .
 \tag{D.54.3}
\]
Here \(D_{b,c}\) is the actual restriction of
the integral derivation with generator values
\((1-s)b,(1-r)c\). On the \(\mathcal T_2\)
component every element of the full central
module \(C\), including \(Z^0\), is fixed.
The closed subgroup generated by these
automorphisms has the continuous elementary
abelian quotient \(\mathscr D\) of D.53,
detected by its action modulo four. Thus its
dimension is the positive arithmetic number
\(d\) bounded in D.53.8.

These are automorphisms of the identity-marked
joint source. The explicit parameter transport
below lifts every difference
\(\alpha_{2(b,c)}\alpha_0^{-1}\) from D.53
through this construction. The chosen D.50
automorphism \(\alpha_0\) itself is not
asserted to have a cusp-fixed extension.
The exact comparison with that fixed
automorphism is the following. Its fixed two
conjugators lift to \(\mathcal X_{54}\);
their two torsion conjugates define an actual
automorphism \(\widetilde\alpha_0\), with
\[
             \widetilde\alpha_0(t)=z_0t,\qquad z_0\in Z .
 \tag{D.54.4}
\]
The cusp defect class
\[
       \omega_0=[z_0]\in Z_G
           =(Z^0)_{G/A}
           \simeq
              \bigl(\bigwedge\nolimits^2 W\bigr)_G
 \tag{D.54.5}
\]
is independent of these lifts of the fixed
conjugators. Composition with any of the
constructed doubled corrections leaves the
literal \(z_0\) unchanged. In particular
those higher even homology corrections do
not change the invariant boundary obstruction.
No value or vanishing assertion for the
actual arithmetic \(\omega_0\) is made.

Throughout the following reduction the
two D.49 conjugators themselves are retained
on \(\mathcal W_h\). Within this fixed framing,
the complete modulo-two pair space
\(\mathcal P\) of D.53 has a well-defined
actual affine quadratic boundary map
\[
 q_{\alpha_0}:\mathcal P\longrightarrow Z_G,
       \qquad q_{\alpha_0}(0)=\omega_0 .
 \tag{D.54.6}
\]
Its polar map, for \(p=(b,c)\) and
\(p'=(b',c')\), is
\[
 q_{\alpha_0}(p+p')+q_{\alpha_0}(p)
       +q_{\alpha_0}(p')+q_{\alpha_0}(0)
       =[(1+s)b'\wedge s(c+b)] .
 \tag{D.54.7}
\]
The right side is interpreted in \(Z_G\).
It vanishes whenever one input belongs to
the exact kernel \(\mathcal K\) in D.53.5.
The map
\[
 \ell:\mathcal K\longrightarrow Z_G,\qquad
     \ell(k)=q_{\alpha_0}(k)+q_{\alpha_0}(0)
 \tag{D.54.8}
\]
is therefore linear, and there is a specified
affine quadratic descent
\[
 \overline q_{\alpha_0}:\mathscr D
           \longrightarrow Z_G/\ell(\mathcal K).
 \tag{D.54.9}
\]
This removes every doubled integral
ambiguity and every same-modulo-two-action
conjugator ambiguity from the displayed
arithmetic boundary map. Its arithmetic
zero locus remains unevaluated in this work;
invertibility of the corresponding
\(\overline T_0+D_p\) is also retained
as part of the arithmetic calculation.

In the separate literal order-\(1080\)
source, the calculation within the retained
\(120\)-parameter moving family is complete:
\[
 \begin{array}{c|r}
 \text{pair parameters}&120\\
 \text{same-action kernel}&102\\
 \text{actual homology parameters}&18\\
 \text{same-action obstruction image}&4\\
 \text{remaining obstruction coordinates}&9\\
 \text{homology parameters tested}&262\,144\\
 \text{full class-two homology images}&864 .
 \end{array}
 \tag{D.54.10}
\]
All \(864\) images are invertible and form
an exactly verified group on homology.
They are exactly the points preserving
the retained modulo-two intersection form.
For every point the certificate gives a
full \(120\)-coordinate conjugator parameter
solving all \(13\) original invariant
equations. The comparison is of zero loci:
the two nine-coordinate polynomial maps
do not have identical coefficient spans.
This evaluated finite classification is
not a numerical evaluation of
\(\overline q_{\alpha_0}\) for \(G_3\).

The source kernel in (D.54.2) remains
nonzero under the explicit next map
\[
 \mathcal T_8=
 \widehat B/
   \overline{\Gamma_3(K_3)(K_3')^4K_3^8}.
 \tag{D.54.11}
\]
For the actual elements \(z_i\in\mathfrak R_h\)
and \(x_j\in K_3\) constructed in D.50,
the elements \([z_i,x_j]^2\) belong to
\(R_{50}\cap L_2\), and their images include
the nonzero module
\[
        2\bigl((U^-/4U^-)\wedge(U/4U)\bigr)
       \ \subset\ \bigwedge\nolimits^2_{\mathbb Z/4}(U/4U).
 \tag{D.54.12}
\]
Full profinite realization remains unresolved.
The manuscript passages are Part D(e),
section 51.16, folio 734 verso, scan 174,
and section 51.15, folios 725--727,
scans 156--160.
\end{statement}
\begin{proof}
\emph{The full joint source and its actual kernels.}
The equality
\(\operatorname{im}(R_{50}\to\mathcal T_2)=Z\)
was proved by the actual preimage and
fourth-power argument in D.51.
The product of the two source kernels
therefore gives the maximal common
marked quotient \(\mathcal T_2/Z\).
Closed Goursat gives (D.54.2), including
both projection kernels. Replacing
\(\mathcal T_2\) by \(\mathcal T_2/Z^0\)
gives exactly D.51's source, so that
projection has kernel \(Z^0\).

\emph{The square module and its moving part.}
The subgroup
\[
                     P_D=D^2D'=D^2
 \tag{D.54.13}
\]
is central in \(D\), abelian of exponent
two, and hence a module over
\(\mathbb F_2[G]\).
Indeed \([u^2,v]=[u,v]^2=1\), and
\((uv)^2u^{-2}v^{-2}\) is a commutator.
Collected class-two coordinates identify
its exact sequence as
\[
 0\longrightarrow\bigwedge\nolimits^2V
   \longrightarrow P_D
   \xrightarrow{\ \pi\ }V\longrightarrow0 .
 \tag{D.54.14}
\]
The last map sends an even abelian
coordinate \(2v\) modulo four to \(v\)
modulo two. Squares of free generators
give every basis vector; its kernel
is exactly the commutator module.
This description includes the square
cross terms and does not choose an
equivariant splitting of the sequence.

Apply the actual central idempotent
\(1-e_A\). Averaging is exact because
\(|A|\) is odd, so
\[
 0\longrightarrow Z^-\longrightarrow
       P_D^-\longrightarrow W\longrightarrow0,
       \qquad P_D^-=(1-e_A)P_D .
 \tag{D.54.15}
\]
For the kernel, decompose
\(V=V^A\oplus W\). The moving part of
\(\bigwedge^2V\) is exactly \(Z^-\):
the summand \(\bigwedge^2(V^A)\) is
fixed, and every other moving summand
lies in \(W\wedge V\).

Let \(E\) denote the image of
\(\mathfrak R_h\) in \(\mathcal T_2\).
D.51 proved \(E\cap C=Z\) and
\(E/Z\simeq U^-/4U^-\) by its actual
abelianization map. Since
\[
         P_D^-=[P_D,A]\subset[D,A]\subset E ,
 \tag{D.54.16}
\]
the image of \(P_D^-\) in
\(\mathcal T_2/Z\) is precisely
the even subgroup \(2(U^-/4U^-)\).
The map is the one in (D.54.14).
This proves the needed agreement with
the D.50 homology coordinate.

\emph{The exact compatible correction module.}
Inside \(\mathcal X_{54}\), form the
actual fibre product
\[
 \mathscr E=
    2U^-\times_W P_D^-,
 \qquad
  2U^-\longrightarrow W,\quad
       2u\longmapsto u\bmod2 .
 \tag{D.54.17}
\]
Its \(\mathcal W_h\) coordinate is
identity. The two remaining coordinates
agree in \(\mathcal T_2/Z\) by
(D.54.16), so this is an embedded
subgroup of the displayed joint source.
It is abelian, centralizes the image
of \(K_3\), and has its \(G\)-action
by the literal conjugation maps.
Its exact sequence is
\[
        0\longrightarrow Z^-\longrightarrow
        \mathscr E\longrightarrow2U^-
        \longrightarrow0 .
 \tag{D.54.18}
\]

Choose the integral pair \((b,c)\)
in D.53.10 and lift \(2b,2c\) to
elements \(B_0,A_0\) of \(\mathscr E\).
Their cusp defect under conjugation
of \(s,r\) is exactly
\[
               (1-s)B_0+s(1-r)A_0 .
 \tag{D.54.19}
\]
There are no unrecorded quadratic
terms: \(\mathscr E\) is abelian and
centralizes the kernel.
The projection of this error in
\(2U^-\) is zero. It therefore lies
in \(Z^-\).
The map
\((1-s)Z^-+s(1-r)Z^-\) is onto
\(Z^-\), since \(Z^-_A=0\).
Correct \(B_0,A_0\) by elements of
this exact kernel to kill (D.54.19).

The resulting conjugates define an
endomorphism of the homology component:
its kernel matrix is \(1+2D_{b,c}\),
which is invertible over \(\mathbb Z_2\).
They also define an endomorphism of
\(\mathcal T_2\), because the image
of \(K_3\) is still class two with
the stated exponent relations.
Its induced map on \(D/\Phi(D)=V\)
is identity. The finite two-group
Frattini criterion makes its restriction
to \(D\) an automorphism; it is identity
on the finite quotient \(G\).
Thus it is an automorphism of
\(\mathcal T_2\). The two component
automorphisms agree on the actual
common quotient through the displayed
conjugators. They give the claimed
automorphism of \(\mathcal X_{54}\).

Its effect on \(C\) is the exterior
square of its effect on \(V\), hence
identity. The parameter has trivial
\(\mathcal W_h\) coordinate.
Its preimage lies in \(\mathfrak R_h
\subset\widehat F'\); equivalently
the reconstruction from the two
conjugators retains the full
normalized commutator condition.
All loop, Fermat, CM and congruence
maps retained by \(\mathcal W_h\)
therefore have identity action.

Modulo four the homology matrices are
\(1+2\overline D_{b,c}\).
Their divided differences add under
composition. D.53.20 proves that the
integral pairs reduce onto all of
\(\mathcal P\), and D.53.4 identifies
their restriction image as \(\mathscr D\).
The generated closed subgroup thus
maps continuously onto its additive
elementary abelian group. The kernel
of this detected quotient is exactly
the identity action on \(U/4U\).
This proves the actual integrated
claim without extending \(\alpha_0\).

\emph{The exact transport of the D.53 differences.}
Write \(\beta_0,A_0\) for the fixed
D.50 conjugators on \(\mathcal Z_h\), and set
\[
 \begin{gathered}
 v_s=\alpha_0^{-1}(\beta_0^{-1}),\qquad
 v_r=\alpha_0^{-1}(A_0^{-1}),\\
 b'=\delta_{b,c}(v_s)+\bar v_s b,\qquad
 c'=\delta_{b,c}(v_r)+\bar v_r c .
 \end{gathered}
 \tag{D.54.23}
\]
Bars here are the actual \(G\)-images,
acting on \(U^-\).
The integral derivation \(\delta_{b,c}\)
factors through \(\mathcal J_\infty\),
since its restriction to \(K_3\) is
abelian and pro-two. Hence it is
defined on these elements of \(\mathcal Z_h\).
The two inverse generator images are
\(v_s s v_s^{-1}\) and \(v_r r v_r^{-1}\).
The crossed-homomorphism identity
therefore gives
\[
 (\delta_{b,c}\circ\alpha_0^{-1})(s)=(1-s)b',
 \qquad
 (\delta_{b,c}\circ\alpha_0^{-1})(r)=(1-r)c'.
\]
Its cusp value is zero because
\(\alpha_0^{-1}(t)=t\).
Thus \((b',c')\in\mathscr P\).
On the complete source \(\mathcal Z_h\),
the D.53 difference is exactly
\(g\mapsto
  2(\delta_{b,c}\circ\alpha_0^{-1})(g)\,g\),
with additive notation in the kernel.
It is the identity-source correction
for \((b',c')\), which has just been
lifted to \(\mathcal X_{54}\).
Its homology matrix is the D.53
matrix \(1+2D_{b,c}T_0^{-1}\).
This proves the stated transport
as an equality of actual maps.

\emph{The fixed arithmetic cusp defect.}
Lift the two specified D.50 conjugators
through \(\mathcal X_{54}\to\mathcal Z_h\).
Their torsion conjugates give a homomorphism
from \(\widehat B\). Its homology component
is the original \(\alpha_0\), and its
\(\mathcal T_2\) component descends
through the characteristic class-two
exponent-four relations of \(K_3\).
Its modulo-two homology map is the
invertible matrix of \(\alpha_0\).
The Frattini criterion again proves
that this component is an automorphism.
The components agree on the common
quotient by construction, giving
\(\widetilde\alpha_0\) on \(\mathcal X_{54}\).
Since \(\alpha_0\) fixes the cusp,
its error belongs to the actual
projection kernel \(Z\), proving
(D.54.4).

Changing either lifted conjugator by
an element of \(Z\) changes its cusp
error by
\((1-s)z_\beta+s(1-r)z_A\).
The image is the augmentation submodule
of \(Z\), so its cokernel is \(Z_G\).
This proves independence in (D.54.5).
Odd averaging kills \(Z^-\) in these
coinvariants. The cross summand
\(W\otimes V^A\) also has zero
\(A\)-coinvariants, which proves the
last isomorphism in (D.54.5).
Finally every constructed doubled
automorphism fixes \(t\) and all of
\(Z\) pointwise. It therefore sends
\(z_0t\) to the same \(z_0t\).
Composition on the other side also
leaves that value unchanged because
the doubled automorphism fixes \(t\).

\emph{The reduced actual arithmetic quadratic map.}
For a pair in \(\mathcal P\), choose
an integral lift in \(\mathscr P\),
which exists by D.53.20.
Use it to change the two D.50
conjugators in the kernel \(U^-\).
These give a cusp-fixed peripheral
endomorphism on \(\mathcal Z_h\);
invertibility is not needed to
define its boundary error.
Lift its conjugators to
\(\mathcal X_{54}\) and take the
class of the cusp error in \(Z_G\).

Changing those lifts contributes
only the augmentation image of \(Z\).
Two integral lifts differ by a
doubled pair. The construction
(D.54.17)--(D.54.19) gives its exact
full class-two correction. Since the
two base actions on \(G\) are identity,
the same calculation applies to
left corrections of the chosen
endomorphism: their action on
\(\mathscr E\) is exactly the
displayed \(s,r,t\) action. It
contributes zero in \(Z_G\).
Thus the error defines the actual
map \(q_{\alpha_0}\) on \(\mathcal P\),
independently of all these choices.

For completeness its polar formula
is the class-two collected product
of the two correction pairs.
Write \(A,B\) for the corrections
of the \(r,s\) conjugators, and
write \(a,b\) for their abelian
coordinates modulo two.
The correction cusp is
\[
                 B\,s(B)^{-1}s(A)t(A)^{-1}.
 \tag{D.54.20}
\]
Replacing \(A,B\) by \(AA',BB'\),
collecting the central commutators,
and using the two zero abelian
cusp values gives the cross term
\((1+s)b'\wedge s(a+b)\).
The actions of the original lifted
conjugates have the same actions
on \(V\) and \(C\) as \(s,r,t\);
their additional terms are linear.
Reducing in \(Z_G\) gives (D.54.7).
This calculation includes the
square terms of the exponent-four
kernel; they are part of the
linear coefficients over \(\mathbb F_2\).

Now let \(p'\in\mathcal K\).
D.53.5 writes it as
\((u_s+d,u_r+d)\) with
\(u_s\in W^s,u_r\in W^r,d\in W^t\).
Its cross term reduces to
\([(1+s)d\wedge s(c+b)]\).
Coinvariance under \(s\) identifies
this with \([d\wedge(1+s)(c+b)]\).
The defining cusp equation for \(p\)
gives
\[
                 (1+s)(c+b)=(1+t)c .
 \tag{D.54.21}
\]
Since \(d\) is fixed by \(t\), its
wedge with this last vector vanishes
in the coinvariants. The polar map
is symmetric, so it vanishes when
either input lies in \(\mathcal K\).
Consequently its restriction to
\(\mathcal K\), after subtracting
the constant value, is linear.
Quotienting its image gives exactly
(D.54.9). This is the actual reduction
of the remaining arithmetic equation;
no zero of that unevaluated map is
asserted.

\emph{The complete separate finite calculation.}
The program
\texttt{d\_cm\_doubled\_class2.py}
first verifies the square-module construction
on this same literal finite source.
It checks all \(543\) actual conjugations
of free-generator squares by \(s,r,z\).
This verifies the full action, including
the central square cross terms, against
the marked \(B\)-words.
The square module has dimension
\(16\,471\), and its moving part has
dimension \(10\,980\).
All \(120\) basis pairs are lifted
through the exact moving square module;
every abelian projection and full cusp
module equation is checked.
An explicit nonzero combination has
divided homology action of rank \(16\)
modulo four. Its entire
\(4\,031\,900\)-letter cusp word and
both torsion relations are evaluated
directly in the full class-two group.
This finite verification is separate
from the symbolic arithmetic proof above.

The program
\texttt{d\_cm\_reduced\_invariant\_equations.py}
reconstructs the literal marked
order-\(1080\) graph. It verifies
the \(102\)-dimensional kernel and
chooses an explicit \(18\)-dimensional
complement in the original
\(120\)-dimensional pair space.
The full retained D.52 quadratic
coefficients are then reduced by
the four-dimensional same-action
obstruction image proved in D.53.
This gives all linear and quadratic
coefficients of the remaining
nine-coordinate map.

For the same eighteen basis parameters,
the program computes every coefficient
of \(T^{\rm t}\Omega T-\Omega\), where
\(\Omega\) is the retained modulo-two
intersection form. Its coefficient
rank is also nine.
An exact comparison of the coefficient
relations records that the two
polynomial spans differ; no equality
of those displayed maps is used.
The separate program
\texttt{d\_cm\_invariant\_zero\_loci.py}
evaluates both maps on all \(2^{18}\)
parameters. Their zero sets coincide
and have exactly \(864\) elements.

For every such parameter the first
program solves the exact linear
equation in the \(102\)-dimensional
same-action kernel. It records the
resulting full \(120\)-coordinate
parameter and checks all original
thirteen invariant equations.
The D.52 central correction construction
then gives full class-two cusp-fixed
conjugators for every point.
Every one of the \(864\) full homology
matrices is explicitly checked to
have rank \(181\).
The finite two-group Frattini criterion
proves generation in the full
class-two quotient.

Finally the program multiplies the
actual homology matrices to build
the subgroup generated by the five
recorded elements
\[
 \mathtt{0x368},\quad\mathtt{0x37e},\quad
 \mathtt{0x3c9},\quad\mathtt{0x635},\quad
 \mathtt{0x80a}.
 \tag{D.54.22}
\]
Each product is reduced in the exact
eighteen-dimensional matrix space
and is checked to belong to the
listed solution set. The generated
set is the full set of \(864\)
invertible matrices. This proves
the asserted group statement.
These calculations concern their
literal finite source throughout.

\emph{A nonzero remaining source-kernel map.}
Choose the actual \(z_i,x_j\) of
D.50's final commutator construction.
The commutator \([z_i,x_j]\) belongs
to \(R_{50}\), and maps to
\(\bar u_i\wedge\bar v_j\) in the
commutator module of \(\mathcal T_2\).
Its square is therefore killed in
\(\mathcal T_2\), so belongs to
\(R_{50}\cap L_2\).
The collected class-two coordinates
in \(\mathcal T_8\) have commutator
module
\(\bigwedge^2_{\mathbb Z/4}(U/4U)\).
The same square maps there to
\(2\bar u_i\wedge\bar v_j\).
Because \(U^-\) is a nonzero integral
direct summand and \(U\) has rank
greater than one, these images
include a nonzero basis coordinate.
Their span is precisely the displayed
submodule in (D.54.12). This proves
the explicit nontriviality claim
without asserting an equality with
the entire remaining image.
\end{proof}
